% Lösungen Einheit 3 von 3 – Symmetrieebenen von Körpern (zusammenbau v0.1)
\einheitenkopf{Einheit 3 von 3 · Symmetrieebenen von Körpern}

\erg{15}{a) $x_2$ – symmetrisch zur $x_1x_3$-Ebene \quad b) Die Paare $A, B$; $C, D$; $E, F$ unterscheiden sich nur im Vorzeichen von $x_2$ – jede Ecke wird auf eine Ecke gespiegelt. \quad c) $M_2$: sie enthält die Diagonale $BD$ und die Spitze $S$ und steht senkrecht auf der Grundfläche. $M_3$, die waagerechte Ebene durch die Giebelspitzen, ist keine: sie spiegelt $S$ auf $(5 | 5 | -1)$. \quad d) $x_1 = 2$; Rechteck $(2 | 0 | 0)$, $(2 | 2 | 0)$, $(2 | 2 | 3)$, $(2 | 0 | 3)$ \quad e) $x_1 = 4$ (Mitte von $1$ und $7$) und $x_2 = 5$ (Mitte von $3$ und $7$) \quad f) $ABC$ ist gleichschenklig mit der Symmetrieachse durch $C$ und die Mitte $(3 | 0 | 0)$ von $AB$; $x_1 = 3$ enthält diese Achse und steht senkrecht auf der Grundfläche – beim geraden Prisma also Symmetrieebene. \quad g) $x_1 = 1$: $A_t, B_t$ und $D, C$ liegen gleich weit links und rechts von $x_1 = 1$, $S_t$ hat $x_1 = 1$. $x_2 = t + 2$: $D, A_t$ und $C, B_t$ haben den Mittelwert $\tfrac{(5 + 2t) + (-1)}{2} = t + 2$, $S_t$ liegt darauf.}
\erg{16}{a) Fehler: die waagerechte Ebene auf halber Höhe ist keine Symmetrieebene – sie spiegelt $S$ auf $(0 | 0 | 0)$, keinen Eckpunkt. Richtig ist etwa $x_1 = 0$ oder $x_1 - x_2 = 0$. \quad b) Eine Symmetrieebene muss jede Ecke auf eine Ecke spiegeln; fällt schon eine Ecke auf keinen Eckpunkt, ist die Bedingung verletzt – ein Gegenbeispiel reicht.}
